\documentclass[11pt]{article}

\usepackage[margin=1in]{geometry}
\usepackage{graphicx}
\usepackage{booktabs}
\usepackage{hyperref}
\usepackage{amsmath}

\newcommand{\maybegraphic}[2]{%
    \IfFileExists{#1}{\includegraphics[width=#2\textwidth]{#1}}{\fbox{Missing file: #1}}%
}

\title{Telegram Ads Forecasting -- Model v1 Report}
\author{Karim Khabib}
\date{\today}

\begin{document}
\maketitle
\tableofcontents

\section{Executive Summary}
The goal is to predict ad views (\texttt{VIEWS}) for Telegram channels given only
three input features available at inference time: \texttt{CPM}, \texttt{CHANNEL\_NAME}, and \texttt{DATE}.
The final Model v1 is an explainable CatBoost baseline that combines
channel-level statistics, CPM-based effects, and yearless seasonality.

Public leaderboard score (best run): \textbf{~0.927}.
The main limitation is label noise and a weak global CPM--VIEWS relationship.
Large gains likely require external channel metadata (subscribers, ER, topic).

\section{Problem Statement}
We must predict \texttt{VIEWS} for each row in the test set.
The training set (\texttt{AllData.csv}) has 7 columns,
while the test set includes only 3 input features.
Therefore, the model must be built using only the features
available at inference time.

\section{Data Overview}
\subsection{Train/Test Summary}
\begin{itemize}
    \item Train size: 142,609 rows, 7 columns.
    \item Test size: 318,722 rows, 4 columns (\texttt{VIEWS} missing).
    \item Missing values: none in key fields.
    \item Date range (train): 2024-10-02 to 2025-12-17 (441 days).
    \item Unique channels (train): 35,957.
    \item Unique channels (test): 6,751 (all appear in train).
\end{itemize}

\subsection{Distributions and Heavy Tails}
Both \texttt{VIEWS} and \texttt{CPM} have heavy tails.
This motivates training on log-transformed targets and features.

\begin{figure}[h]
    \centering
    \maybegraphic{figures/views_distribution.png}{0.85}
    \caption{Distribution of VIEWS (linear and log scales).}
\end{figure}

\begin{figure}[h]
    \centering
    \maybegraphic{figures/cpm_distribution.png}{0.85}
    \caption{Distribution of CPM (linear and log scales).}
\end{figure}

\subsection{Channel Sparsity}
Channels are highly sparse: median rows per channel = 2,
90\% have at most 9 rows, and the maximum is 137.
This requires shrinkage toward global statistics to avoid overfitting.

\begin{figure}[h]
    \centering
    \maybegraphic{figures/channel_counts_hist.png}{0.85}
    \caption{Histogram of rows per channel.}
\end{figure}

\section{EDA Findings}
\subsection{Global vs. Within-Channel CPM Relationship}
Global correlation between \texttt{CPM} and \texttt{VIEWS} is weak.
However, within-channel correlations are often positive.
This suggests that CPM is informative only after conditioning on channel.

\begin{figure}[h]
    \centering
    \maybegraphic{figures/corr_heatmap.png}{0.85}
    \caption{Pearson/Spearman correlation on log-transformed variables.}
\end{figure}

\begin{figure}[h]
    \centering
    \maybegraphic{figures/cpm_views_binned.png}{0.85}
    \caption{Binned CPM -- VIEWS dependency (median and IQR).}
\end{figure}

\subsection{Seasonality}
Yearless seasonality (day-of-week and day-of-year cycles) shows mild patterns,
so we keep simple date features without absolute time trends.

\begin{figure}[h]
    \centering
    \maybegraphic{figures/dow_seasonality.png}{0.75}
    \caption{Average VIEWS by day of week.}
\end{figure}

\section{Modeling Approach}
\subsection{Target Transformation}
We predict \(\log(1 + \texttt{VIEWS})\) and invert predictions with \(\exp(x) - 1\).
This stabilizes variance and reduces the impact of extreme outliers.

\subsection{Features}
\begin{itemize}
    \item \textbf{CPM features:} \texttt{log\_cpm}, \texttt{cpm\_to\_ch\_median}.
    \item \textbf{Channel features:} smoothed median of log views,
    smoothed CTR, smoothed actions rate, CPM slope,
    and expected log prediction \texttt{ch\_log\_pred}.
    \item \textbf{Date features:} \texttt{dow}, \texttt{is\_weekend}, \texttt{month},
    \texttt{doy\_sin}, \texttt{doy\_cos}.
\end{itemize}

All channel-level statistics are computed on train only and smoothed
with shrinkage factors to avoid overfitting.

\subsection{Model and Post-Processing}
\begin{itemize}
    \item Model: \texttt{CatBoostRegressor} with MAE objective.
    \item Post-processing: prediction blending with channel baselines,
    optional clipping by per-channel quantiles, optional scaling.
\end{itemize}

\section{Training and Validation}
We use a time-based holdout of the last 30 days for quick local checks.
However, the public test set contains earlier dates, so local validation
is not fully aligned with the leaderboard.

\section{Metric Exploration}
The public metric is not documented. We probed it with constant predictions:

\begin{center}
\begin{tabular}{lr}
\toprule
Constant prediction & Public score \\
\midrule
0 & 1.0000 \\
300 & 0.9427 \\
1000 & 1.0938 \\
3000 & 1.9048 \\
\bottomrule
\end{tabular}
\end{center}

This behavior suggests a relative-error metric (MAPE/SMAPE-like),
where large overestimates are heavily penalized.
Therefore, conservative predictions and clipping help.

\section{Results}
Best public score achieved with Model v1: \textbf{~0.927}.
This is stronger than naive baselines but still far from ideal,
primarily due to noise and missing channel metadata.

\section{Error Analysis}
\begin{itemize}
    \item The same (\texttt{CPM}, \texttt{CHANNEL\_NAME}, \texttt{DATE})
    can map to different \texttt{VIEWS}, indicating label noise.
    \item Most channels have very few examples, so channel statistics are noisy.
    \item CPM is weak globally; it matters mostly within channels.
\end{itemize}

\section{Limitations}
\begin{itemize}
    \item Only three features are available at inference time.
    \item Channel metadata (subscribers, ER, topic) is missing.
    \item Validation is misaligned with the test temporal distribution.
    \item The evaluation metric is unknown and favors conservative predictions.
\end{itemize}

\section{Future Work}
\begin{itemize}
    \item Enrich channels with external metadata (subscribers, ER, topic).
    \item Try monotonic constraints or isotonic regression for CPM within channels.
    \item Use per-channel KNN over CPM as a local smoother.
    \item Build a better validation scheme aligned with test distribution.
\end{itemize}


\end{document}
